\documentclass[12pt,a4paper]{article}

\usepackage[utf8]{inputenc}
\usepackage[T2A]{fontenc}
\usepackage[english,russian]{babel}
\usepackage{array}
\usepackage{booktabs}
\usepackage{caption}
\usepackage{float}
\usepackage{geometry}
\usepackage{hyperref}
\usepackage{indentfirst}
\usepackage{qrcode}
\usepackage{tabularx}
\usepackage{tocloft}
\geometry{margin=25mm}
\hypersetup{unicode=true,colorlinks=true,urlcolor=blue,linkcolor=black}
\setlength{\parindent}{1.25cm}
\setlength{\parskip}{0pt}
\setlength{\emergencystretch}{1em}
\setcounter{tocdepth}{2}
\setlength{\cftbeforesecskip}{5pt}
\setlength{\cftbeforesubsecskip}{1pt}
\renewcommand{\cftsecfont}{\normalsize\bfseries}
\renewcommand{\cftsecpagefont}{\normalsize\normalfont}
\renewcommand{\cftsubsecfont}{\normalsize\normalfont}
\renewcommand{\cftsubsecpagefont}{\normalsize\normalfont}
\renewcommand{\cftdotsep}{1}
\renewcommand{\labelitemi}{\textendash}
\DeclareCaptionLabelSeparator{emdash}{\space---\space}
\captionsetup[figure]{
  justification=centering,
  singlelinecheck=false,
  labelsep=emdash
}
\addto\captionsrussian{\renewcommand{\figurename}{Рис.}}
\newcommand{\code}[1]{\texttt{\detokenize{#1}}}
\newcolumntype{P}[1]{>{\raggedright\arraybackslash}p{#1}}

\begin{document}
\begin{titlepage}
\begin{center}
{\large МОСКОВСКИЙ АВИАЦИОННЫЙ ИНСТИТУТ\\ (НАЦИОНАЛЬНЫЙ ИССЛЕДОВАТЕЛЬСКИЙ УНИВЕРСИТЕТ)}

\vspace{18pt}
{\large Факультет информационных технологий и прикладной математики}

{\large Кафедра вычислительной математики и~программирования}

\vfill
{\large \textbf{Лабораторная работа №2 по курсу\\ «Информационный поиск»}}

\vspace{12mm}
{\large \textbf{Поисковый робот}}

\vfill
\begin{flushright}
\begin{tabular}{@{}r@{:\hspace{0.6em}}l@{}}
\textit{Студент} & К.\,А. Арусланов\\
\textit{Группа} & М8О-403Б-23\\
\textit{Преподаватель} & А.\,А. Кухтичев\\
\textit{Дата} & 08.10.2026\\
\textit{Оценка} &
\end{tabular}
\end{flushright}
\vfill
Москва, 2026
\end{center}
\end{titlepage}

\tableofcontents
\newpage

\section{Постановка задачи}

Для корпуса игровых публикаций из первой лабораторной работы требуется
разработать поисковый робот:

\begin{itemize}
  \item реализовать обкачку документов;
  \item принимать при обычном запуске один аргумент --- путь к YAML-файлу с
  данными базы в секции \code{db} и параметрами робота в секции \code{logic},
  в том числе задержкой между запросами;
  \item сохранять в базе нормализованный URL, сырой HTML, название источника
  и дату загрузки в формате Unix timestamp;
  \item после остановки продолжать обход с незавершённого документа;
  \item периодически проверять загруженные страницы и обновлять сохранённые
  документы только при изменении содержимого.
\end{itemize}

\section{Метод решения}

\subsection{Источники и входные данные}

В качестве корпуса выбраны публикации восьми игровых сайтов: PlayGround,
StopGame, IXBT.games, Игромания, GameMAG, PC Gamer, Eurogamer и
GamingOnLinux. Входом служит UTF-8 TSV-файл: короткий \code{urls.txt} из
первой лабораторной либо объединённый инвентарь sitemap. Последний содержит
1\,055\,674 URL-кандидата, но это ещё не число пригодных документов.

Файл читается построчно и импортируется в MongoDB партиями, без загрузки
всего списка в память. URL очищается от параметров отслеживания, его хост
сверяется с разрешённым списком. Для новостей StopGame дополнительно
исправляется встречающийся в sitemap префикс \code{/news/} на рабочий
\code{/newsdata/}. Пробная выборка воспроизводима по заданному seed;
порядок очереди перемешивается устойчивым хешем URL.

\subsection{Очередь и сохранение документов}

Для хранения выбрана MongoDB. Очередь и содержимое страницы разделены по
коллекциям, чтобы просмотр результатов разбора не загружал большой HTML:

\begin{center}
\small
\renewcommand{\arraystretch}{1.13}
\begin{tabularx}{\textwidth}{@{}P{0.25\textwidth}X@{}}
\toprule
Коллекция & Содержимое \\
\midrule
\code{frontier} & URL, источник, состояние, время следующей попытки,
число попыток, HTTP-заголовки и последняя ошибка \\
\code{raw_pages} & Нормализованный URL, источник, Unix-время загрузки,
сырой HTML и SHA-256 ответа \\
\code{documents} & Результат разбора, HTTP-метаданные, время загрузки и
ссылка \code{raw_page_id} на исходный HTML \\
\code{crawler_meta} & Метка импорта инвентаря и паузы для отдельных сайтов \\
\bottomrule
\end{tabularx}
\end{center}
\normalsize

Ключ \code{_id} у заданий и документов равен нормализованному URL, поэтому
повторная загрузка не создаёт дубликатов. Сначала сохраняется сырой HTML,
затем результат разбора и только после этого завершается задание очереди.
Если выполнение прервётся между этими шагами, при следующей попытке
согласованность записей проверяется по SHA-256.

Робот атомарно выбирает доступный URL и отмечает его как обрабатываемый.
Затем он проверяет \code{robots.txt} и делает запрос с паузой для каждого
хоста. Редиректы проверяются по отдельности: переход за пределы источника
не выполняется. При перезапуске незавершённые URL возвращаются в начало
очереди. Механизм
рассчитан на один процесс робота.

\subsection{Проверка изменений}

При повторном запросе робот использует сохранённые \code{ETag} и
\code{Last-Modified}. Ответ \code{304} оставляет HTML неизменным. Если сервер
возвращает \code{200}, вычисляется SHA-256 исходного ответа: при совпадении
хеша обе записи не перезаписываются. Для изменённого ответа обновляются
\code{raw_pages}, \code{documents} и \code{fetched_at}. Если одна из двух
записей отсутствует либо их хеши расходятся, робот повторно получает HTML
и восстанавливает согласованное состояние. Время каждой проверки хранится в
\code{frontier.last_checked_at}. При ошибках попытка откладывается:
\code{404}/\code{410} и запрет \code{robots.txt} проверяются реже, для
\code{429} учитывается \code{Retry-After}. Пауза одного источника хранится
в MongoDB и не останавливает другие сайты.

\subsection{Проверка работы}

При переносе компонентов из \code{lab1} в \code{lab2} добавлены парсеры
GameMAG, GamingOnLinux, PC Gamer и Eurogamer. Для нового робота написаны
отдельные тесты с имитацией MongoDB и HTTP-ответов. Проверяются:

\begin{itemize}
  \item чтение YAML и оба формата входного TSV;
  \item загрузка и сохранение обязательных полей;
  \item повторный запуск после завершения и после остановки на URL;
  \item ответы \code{304}, неизменённый \code{200} и изменённый \code{200};
  \item запрет \code{robots.txt}, отличение текстовых правил от HTML-заглушки
  и отложенная попытка после \code{429};
  \item пауза для одного источника без остановки других сайтов и
  воспроизводимость выборки URL.
\end{itemize}

\section{Результаты}

\subsection{Запуск и автотесты}

Код робота запускается командой \code{python crawler.py crawler.example.yaml}.
В примере ограничение \code{max_documents: 10} позволяет проверить небольшую
часть списка \code{urls.txt}. Команда \code{python run_tests.py} проверяет
систему без сервера MongoDB и без обращений к сайтам.

Успешно прошли 48 тестов, из них 18 проверяют робота. Сценарии охватывают
сохранение HTML и времени Unix, продолжение после остановки, ответы
\code{304} и \code{200}, контрольную сумму, \code{robots.txt} и
\code{Retry-After}. Локальные HTML-образцы первой работы проверяются
прежними тестами.

\subsection{Пробная обкачка}

MongoDB запускалась локально через Docker Compose. Сначала робот сохранил
24 проверочные страницы из \code{urls.txt}, затем выборка была расширена.
На 8 октября 2026 года в очереди находилось 824 URL: 374 завершённых,
449 ожидающих и один оставшийся в обработке после остановки. Большой
инвентарь в миллион URL в базу пока не импортировался.

\begin{center}
\small
\renewcommand{\arraystretch}{1.12}
\begin{tabular}{@{}lrrr@{}}
\toprule
Источник & Скачано & С текстом & Без текста \\
\midrule
PlayGround & 46 & 42 & 4 \\
StopGame & 51 & 51 & 0 \\
IXBT.games & 46 & 46 & 0 \\
Игромания & 36 & 36 & 0 \\
GameMAG & 54 & 53 & 1 \\
PC Gamer & 56 & 52 & 4 \\
Eurogamer & 39 & 36 & 3 \\
GamingOnLinux & 46 & 45 & 1 \\
\midrule
\textbf{Всего} & \textbf{374} & \textbf{361} & \textbf{13} \\
\bottomrule
\end{tabular}
\end{center}
\normalsize

Все 374 записи \code{documents} имеют соответствующие страницы
\code{raw_pages}; их контрольные суммы совпадают, пустого HTML и
потерянных исходных страниц нет. У 13 документов HTML получен, но парсер
не нашёл основной текст. Среди них есть карточки игр и страницы категорий,
которые нужно исключить из инвентаря, а также публикации с необычной
разметкой или преимущественно медийным содержанием. Поэтому для будущего
поискового индекса пригодность URL надо проверять после загрузки, а не
оценивать только по sitemap.

Первоначальные 124 записи были перенесены из формата со встроенным HTML
в раздельные коллекции после резервного копирования базы. Средний размер
записи в \code{documents} уменьшился примерно с 470 до 7~КБ, что удобнее
для просмотра результатов разбора. Дополнительно 71\,681 URL новостей
StopGame в полном инвентаре и 57 старых заданий очереди приведены к
рабочему виду \code{/newsdata/}; исходные задания сохранены для
восстановления.

\newpage
\subsection{Репозиторий}

Исходный код, конфигурации, тесты и отчёт размещены в ветке \code{lab2}
репозитория: \url{https://github.com/kirbasss/MAI-IR/tree/lab2}.
Сырые страницы и содержимое базы данных в GitHub не загружаются.

\begin{figure}[H]
\centering
\qrcode*[height=28mm]{https://github.com/kirbasss/MAI-IR/tree/lab2}
\caption{QR-код для перехода к ветке \code{lab2} репозитория.}
\end{figure}

\section{Вывод}

В этой работе я дополнил сбор игровых публикаций устойчивой очередью в MongoDB.
Теперь робот может продолжить обход после остановки и периодически проверять
уже сохранённые страницы. Было полезно разобраться, как связать состояние
очереди с сохранением документа и зачем при повторной проверке нужны
\code{ETag}, \code{Last-Modified} и контрольная сумма ответа.

Я проверил робота тестами и пробной загрузкой в MongoDB. На практике
выяснилось, что даже успешный ответ сайта не всегда означает пригодный
текстовый документ: в sitemap встречаются страницы категорий, а у некоторых
публикаций иная разметка. Перед массовой обкачкой нужно уточнить фильтры
URL и правила обработки таких страниц, а также следить за расходом места.

\end{document}
